\begin{example}[Randomized controlled trial]
    The randomized controlled trial is a joint distribution of variables
    \begin{itemize}
        \item Confounding variable $\catvariableof{C}$ summarizing other factors affecting the outcome.
        \item Treatment assignment $\catvariableof{T}$ with values in $\{0,1\}$.
        Here, $1$ indicates assignment to the treatment group and $0$ to the control group.
        The experiment is conducted such that the treatment variable is independent of the confounding variable $\catvariableof{C}$.
        \item Outcome $\catvariableof{O}$ .
    \end{itemize}
    We want to estimate the causal effect of the treatment on the outcome, namely calculate the intervention query
    \begin{align*}
        \condprobat{\catvariableof{O}}{\dovariableof{T}=1} \, .
    \end{align*}
    By design of the randomized control, the confounding variable has no causal effect on $\catvariableof{T}$.
    Thus, we have
    \begin{align*}
        \condprobat{\catvariableof{O}}{\dovariableof{T}=1}
        = \condprobat{\catvariableof{O}}{\catvariableof{T}=1}
    \end{align*}
    that is, the intervention query is equal to the conditional probability (see \figref{fig:RCT}).


    % If confounder influences treatment
    If the treatment variable would depend on the confounder (see \figref{fig:RCTwithConfounder}), the effect of the intervention can not be identified based on conditional probabilities, when not measuring the confounder $\catvariableof{C}$.

    \begin{figure}[hbt!]
        \stantikzpic{

            \drawtextnode{-9}{0}{$a)$}
            \begin{scope}[shift={(-2,-1)}]

                \drawvariablenode{0}{0}{$\catvariableof{O}$}{XO}
                \drawvariablenode{-7}{-5}{$\dovariableof{T}$}{DT}
                \drawvariablenode{-3}{-5}{$\catvariableof{T}$}{XT}
                \drawvariablenode{3}{-5}{$\catvariableof{C}$}{XC}

                \draw[->-] (DT) -- (XT);

                \draw[->-] (XT) -- (0,-2);
                \draw[->-] (0,-2) -- (XO);
                \draw[->-] (XC) -- (0,-2);

            \end{scope}


            \drawtextnode{10}{0}{$b)$}
            \begin{scope}[shift={(17,0)}]
                \drawtdwire{0}{0}{$\catvariableof{0}$}
                \drawtensorblock{0}{-3}{3}{$\condprobat{\catvariableof{O}}{\catvariableof{T},\catvariableof{C}}$}

                \drawtdwire{-2}{-4}{$\catvariableof{T}$}
                \drawtensorblock{-2}{-7}{3}{$\condprobat{\catvariableof{T}}{\dovariableof{T}}$}
                \drawhdwire{-5}{-7}{$\dovariableof{T}$}
                \drawsquaretensor{-8}{-7}{$\onehotmapof{\catindexof{T}}$}

                \drawtdwire{2.5}{-4}{$\catvariableof{C}$}
                \drawtensorblock{3}{-7}{1.5}{$\probat{\catvariableof{C}}$}
%                \drawvariabledot{2}{-6}
            \end{scope}

            \drawtextnode{22}{-5}{${=}$}

            \begin{scope}[shift={(25,0)}]
                \drawtdwire{0}{0}{$\catvariableof{0}$}
                \drawtensorblock{0}{-3}{2.5}{$\condprobat{\catvariableof{O}}{\catvariableof{T}}$}
                \drawtdwire{0}{-4}{$\catvariableof{T}$}
                \drawsquaretensor{0}{-7}{$\onehotmapof{\catindexof{T}}$}
                %\drawtdwire{2}{-4}{$\catvariableof{C}$}
                %\drawvariabledot{2}{-6}
            \end{scope}

        }
        \caption{Model of a Randomized Controlled Trial:
        The confounding variable $\catvariableof{C}$ summarizes all influenced on the outcome $\catvariableof{O}$ other than the treatment variable $\catvariableof{T}$.
        a)
            b) Effect of intervening on the treatement variable $\catindexof{T}\in\{0,1\}$}\label{fig:RCT}
    \end{figure}



    \begin{figure}[hbt!]
        \stantikzpic{

            \drawtextnode{-9}{0}{$a)$}
            \begin{scope}[shift={(-2,-1)}]

                \drawvariablenode{0}{0}{$\catvariableof{O}$}{XO}
                \drawvariablenode{-7}{-6}{$\dovariableof{T}$}{DT}
                \drawvariablenode{-3}{-3}{$\catvariableof{T}$}{XT}
                \drawvariablenode{1}{-6}{$\catvariableof{C}$}{XC}

                \draw[->-] (DT) -- (-3,-5);
                \draw[->-] (-3,-5) -- (XT);
                \draw[->-] (XC) -- (-3,-5);

                \draw[->-] (XT) -- (0,-2);
                \draw[->-] (0,-2) -- (XO);
                \draw[->-] (XC) -- (0,-2);

            \end{scope}


            \drawtextnode{10}{0}{$b)$}
            \begin{scope}[shift={(17,0)}]
                \drawtdwire{0}{0}{$\catvariableof{0}$}
                \drawtensorblock{0}{-3}{3}{$\condprobat{\catvariableof{O}}{\catvariableof{T},\catvariableof{C}}$}

                \drawtdwire{-2}{-4}{$\catvariableof{T}$}
                \drawtensorblock{-2}{-7}{3}{$\condprobat{\catvariableof{T}}{\catvariableof{C},\dovariableof{T}}$}
                \drawhdwire{-5}{-7}{$\dovariableof{T}$}
                \drawsquaretensor{-8}{-7}{$\onehotmapof{\catindexof{T}}$}

                \drawtdwire{2}{-4}{$\catvariableof{C}$}
                \draw  (2,-10) -- (2,-6);
                \drawvariabledot{2}{-10}
                \draw [->-] (2,-11) -- (2,-10);
                \draw [->-] (2,-10) to [bend right=-30] (-2,-8);
                \drawtensorblock{2}{-12}{1.5}{$\probat{\catvariableof{C}}$}

            \end{scope}

            \drawtextnode{22}{-5}{${=}$}

            \begin{scope}[shift={(27,0)}]
                \drawtdwire{0}{0}{$\catvariableof{0}$}
                \drawtensorblock{0}{-3}{3}{$\condprobat{\catvariableof{O}}{\catvariableof{T},\catvariableof{C}}$}
                \drawtdwire{-2}{-4}{$\catvariableof{T}$}
                \drawsquaretensor{-2}{-7}{$\onehotmapof{\catindexof{T}}$}
                \drawtdwire{2}{-4}{$\catvariableof{C}$}
                %\drawvariabledot{2}{-6}
                %\drawtdwire{2.5}{-4}{$\catvariableof{C}$}
                \drawtensorblock{2}{-7}{1.5}{$\probat{\catvariableof{C}}$}
            \end{scope}

        }
        \caption{Effect when the confounder has a causal influence on the treatment:
        The causal effect is not identifiable when not measuring the confounding variable.}\label{fig:RCTwithConfounder}
    \end{figure}

\end{example}